\documentclass[11pt]{article}
\usepackage[a4paper,margin=1in]{geometry}
\usepackage{amsmath,amssymb,amsthm,microtype}
\usepackage{booktabs}
\usepackage{cleveref}
\newtheorem{theorem}{Theorem}[section]
\newtheorem{corollary}[theorem]{Corollary}
\newtheorem{proposition}[theorem]{Proposition}
\newtheorem{lemma}[theorem]{Lemma}
\theoremstyle{definition}
\newtheorem{definition}[theorem]{Definition}
\theoremstyle{remark}
\newtheorem{remark}[theorem]{Remark}
\newcommand{\R}{\mathbb R}
\newcommand{\1}{\mathbf 1}
\newcommand{\cN}{\mathcal N}
\title{Sharp Physical Repelling-Return Affine-Leg Remainder}
\author{Bin Wang}
\date{August 2026}
\begin{document}
\maketitle

\begin{abstract}
The first physical endpoint leg of RH-324 produces an actual finite-noise
prefix near the repelling point.  We pass that same prefix through the
second physical row and compare only this row with the RH-323 repelling
tangent row, retaining the incoming coordinate.  After including the
Jacobian, the scaled physical row has an exact two-lobe formula.  Its
$L^1$ distance to the full Gaussian at the curved image is exactly two
state-boundary tails, while removal of the quadratic displacement costs
\(4\Phi(u_c\sigma u^2/2)-2\).  We prove a uniform fourth moment for the
actual first-leg marginal; this is the missing uniform-integrability input
that ordinary $L^1$ convergence does not provide.  For either repelling
orientation,
\[
 \frac{D_{\sigma,a_\sigma}^{\pm}}{\sigma}
 \longrightarrow
 \sqrt{\frac2\pi}\,u_c
 \int_{\{\pm u>0\}}u^2p_d(u)\,du>0,
 \qquad a_\sigma\longrightarrow d.
\]
The sum has a closed coefficient.  Thus exponential, or even
$o(\sigma)$, accuracy is false for this exact second hybrid Duhamel row
term.  The coefficient vanishes at the single row $u=0$, and a source
moving to the critical partition disproves a global uniform $O(\sigma)$
row estimate.  No fully physical-to-fully affine two-leg equality,
all-cycle transport, cyclic trace estimate, or Gate A--E promotion is
asserted.
\end{abstract}

\section{Physical row and repelling coordinates}

Let $u_c$ be the root in $(1,2)$ of
\[
 u^3-2u^2+2u-2=0,
\]
and put
\begin{equation}\label{eq:constants}
 f(x)=1-u_cx^2,
 \qquad r=u_c-1,
 \qquad \lambda=2u_cr,
 \qquad \alpha=2u_c,
 \qquad b=u_c^{-1/2}.
\end{equation}
Numerically,
\begin{align*}
 u_c&=1.543689012692076\ldots,
 &r&=0.543689012692076\ldots,\\
 \lambda&=1.678573510428322\ldots,
 &b&=0.804859535112210\ldots .
\end{align*}
Write $\phi$, $\Phi$, and $\overline\Phi=1-\Phi$ for the standard normal
density, distribution, and survival functions.  The archived row-normalized
folded kernel on $[0,1]$ is
\begin{equation}\label{eq:physical-kernel}
 P_\sigma(x,y)=
 \frac{\phi_\sigma(y-f(x))+\phi_\sigma(y+f(x))}{Z_\sigma(x)},
 \qquad
 \phi_\sigma(t)=\sigma^{-1}\phi(t/\sigma),
\end{equation}
where
\begin{equation}\label{eq:normalizer}
 Z_\sigma(x)=
 1-\overline\Phi\!\left(\frac{1-f(x)}\sigma\right)
  -\overline\Phi\!\left(\frac{1+f(x)}\sigma\right).
\end{equation}
This is the probability that the folded Gaussian lies in the physical
state interval; in particular $0<Z_\sigma(x)\le1$.

Use repelling coordinates
\begin{equation}\label{eq:repelling-coordinates}
 x=r+\sigma u,
 \qquad y=r+\sigma w,
 \qquad
 I_\sigma=\left[-\frac r\sigma,\frac{1-r}\sigma\right].
\end{equation}
The quadratic map gives the exact, rather than asymptotic, expansion
\begin{equation}\label{eq:repelling-expansion}
 f(r+\sigma u)=r-\lambda\sigma u-u_c\sigma^2u^2
 =r-\sigma d_\sigma(u),
 \qquad
 d_\sigma(u)=\lambda u+u_c\sigma u^2.
\end{equation}
The Jacobian in the next definition is essential.

\begin{definition}[Second physical, curved, and tangent rows]
For $u\in I_\sigma$, define the scaled physical row
\begin{align}
 L_\sigma(u,w)
 &:=\sigma P_\sigma(r+\sigma u,r+\sigma w)\nonumber\\
 &=\frac{
 \phi(w+d_\sigma(u))
 +\phi\!\left(w+\frac{2r}\sigma-d_\sigma(u)\right)}
 {Z_\sigma(r+\sigma u)}\1_{I_\sigma}(w),
 \label{eq:physical-second-row}
\end{align}
and, on the full line,
\begin{equation}\label{eq:curved-tangent-rows}
 C_{\sigma,u}(w)=\phi(w+d_\sigma(u)),
 \qquad
 A_u(w)=\phi(w+\lambda u).
\end{equation}
All three rows have mass one.
\end{definition}

The first lobe in \eqref{eq:physical-second-row} is centered at the signed
exact image; that image is positive in the fixed local repelling chart near
$r$.  The other lobe is the folded companion.  No symmetry with the endpoint
expansion of RH-324 has been assumed.

\section{Exact row identities and their sharp local limit}

\begin{theorem}[Exact repelling-row identities]\label{thm:row-identities}
For every $\sigma>0$ and every physical source $u\in I_\sigma$,
\begin{equation}\label{eq:boundary-identity}
 \boxed{
 \|L_\sigma(u,\cdot)-C_{\sigma,u}\|_1
 =2\left[
 \overline\Phi\!\left(\frac r\sigma-d_\sigma(u)\right)
 +\overline\Phi\!\left(\frac{1-r}\sigma+d_\sigma(u)\right)
 \right].}
\end{equation}
Moreover,
\begin{equation}\label{eq:shift-identity}
 \boxed{
 \|C_{\sigma,u}-A_u\|_1
 =4\Phi\!\left(\frac{u_c\sigma u^2}{2}\right)-2.}
\end{equation}
Consequently, if $B_\sigma(u)$ and $S_\sigma(u)$ denote the right sides of
\eqref{eq:boundary-identity} and \eqref{eq:shift-identity}, then
\begin{equation}\label{eq:triangle-sandwich}
 \max\{0,S_\sigma(u)-B_\sigma(u)\}
 \le \|L_\sigma(u,\cdot)-A_u\|_1
 \le S_\sigma(u)+B_\sigma(u).
\end{equation}
\end{theorem}

\begin{proof}
The two unnormalized lobes in \eqref{eq:physical-kernel} have total mass one
on $[0,\infty)$, so restriction to $[0,1]$ gives $Z_\sigma\le1$.
On $I_\sigma$, the numerator of $L_\sigma$ contains
$C_{\sigma,u}$ and division by $Z_\sigma\le1$ only increases it.  Hence
$L_\sigma\ge C_{\sigma,u}$ on $I_\sigma$, while $L_\sigma=0$ off that
interval.  Since both are probability densities, the positive and negative
parts of their difference have equal mass.  Each mass is the amount of
$C_{\sigma,u}$ outside $I_\sigma$, namely
\[
 \overline\Phi(r/\sigma-d_\sigma(u))
 +\overline\Phi((1-r)/\sigma+d_\sigma(u)).
\]
This proves \eqref{eq:boundary-identity}; in particular the companion lobe,
state restriction, and exact normalizer have already been combined.

The centers in \eqref{eq:curved-tangent-rows} differ by
$u_c\sigma u^2$.  Unit Gaussians separated by $t$ have unhalved $L^1$
distance $4\Phi(|t|/2)-2$, which proves \eqref{eq:shift-identity}.
The ordinary and reverse triangle inequalities give
\eqref{eq:triangle-sandwich}.
\end{proof}

\begin{corollary}[Sharp fixed-row coefficient]\label{cor:fixed-row}
For each fixed $u\in\R$,
\begin{equation}\label{eq:fixed-row-limit}
 \boxed{
 \lim_{\sigma\downarrow0}
 \frac{\|L_\sigma(u,\cdot)-A_u\|_1}{\sigma}
 =\sqrt{\frac2\pi}\,u_cu^2.}
\end{equation}
At $u=0$ the coefficient is zero; for every fixed $u\ne0$ it is positive.
\end{corollary}

\begin{proof}
For fixed $u$, both arguments of the survival functions in
\eqref{eq:boundary-identity} are a positive constant times $\sigma^{-1}$
up to $O(1)$.  Thus $B_\sigma(u)=o(\sigma)$.  Taylor expansion at zero gives
\[
 \frac{S_\sigma(u)}\sigma
 \longrightarrow 2\phi(0)u_cu^2
 =\sqrt{2/\pi}\,u_cu^2.
\]
Now use \eqref{eq:triangle-sandwich}.
\end{proof}

The fixed-source quantifier cannot be made global.

\begin{proposition}[No global uniform $O(\sigma)$ row theorem]
\label{prop:global-obstruction}
Let
\[
 u_\sigma=\frac{b-r}{\sigma},
\]
so that the physical source $r+\sigma u_\sigma$ equals the critical
partition $b$.  Then
\begin{equation}\label{eq:global-obstruction}
 \liminf_{\sigma\downarrow0}
 \|L_\sigma(u_\sigma,\cdot)-A_{u_\sigma}\|_1\ge1.
\end{equation}
In particular,
$\sup_{u\in I_\sigma}\|L_\sigma(u,\cdot)-A_u\|_1$ is not
$O(\sigma)$.
\end{proposition}

\begin{proof}
Since $f(b)=0$, \eqref{eq:repelling-expansion} gives
$d_\sigma(u_\sigma)=r/\sigma$.  Therefore
$B_\sigma(u_\sigma)\to1$: the curved Gaussian has half its mass beyond the
state boundary, while the other boundary tail vanishes.  On the other hand,
the curved-to-tangent displacement is
\[
 u_c\sigma u_\sigma^2=\frac{u_c(b-r)^2}{\sigma}\longrightarrow\infty,
\]
so $S_\sigma(u_\sigma)\to2$.  The lower bound in
\eqref{eq:triangle-sandwich} yields \eqref{eq:global-obstruction}.
\end{proof}

\section{The actual physical first-leg prefix}

We now specify the incoming law rather than reusing a local seed.  Put
$L=\sigma^{-1}$ and, for $0\le a\le L$, let the exact RH-322 entrance
density be
\begin{equation}\label{eq:finite-seed}
 \widetilde h_{\sigma,a}(v)=
 \frac{\phi(v-a)+\phi(2L-v-a)}
 {\Phi(a)-\overline\Phi(2L-a)}\1_{[0,L]}(v).
\end{equation}
For $c_\sigma(v)=\alpha v-u_c\sigma v^2$, the exact RH-324 first physical
row is
\begin{equation}\label{eq:first-row}
 K_\sigma(v,u)=
 \frac{
 \phi(u+2r/\sigma-c_\sigma(v))+\phi(u+c_\sigma(v))}
 {Z_\sigma(1-\sigma v)}\1_{I_\sigma}(u).
\end{equation}
Define its actual $U$-marginal by
\begin{equation}\label{eq:actual-prefix}
 \mu_{\sigma,a}(du)=
 \left\{\int_0^L\widetilde h_{\sigma,a}(v)
 K_\sigma(v,u)\,dv\right\}du.
\end{equation}
This is an exact physical prefix, not the affine $U$-law.

For $d\ge0$, the RH-323 affine entrance and intermediate densities are
\begin{align}
 g_d(v)&=\frac{\phi(v-d)}{\Phi(d)}\1_{[0,\infty)}(v),
 \label{eq:gd}\\
 p_d(u)&=\frac1{s_1\Phi(d)}
 \phi\!\left(\frac{u+\alpha d}{s_1}\right)
 \Phi\!\left(\frac{d-\alpha u}{s_1}\right),
 \qquad s_1=\sqrt{1+\alpha^2}.
 \label{eq:pd}
\end{align}
Equivalently, $p_d$ is the law of
\begin{equation}\label{eq:affine-U}
 U=-\alpha V-Z_1,
 \qquad V\sim g_d,
 \qquad Z_1\sim N(0,1),
\end{equation}
with $V$ and $Z_1$ independent.  RH-324 proves, for fixed
$0<\rho<1-b$ and
\[
 m_\rho=u_c(1-\rho)^2-1>0,
 \qquad
 M_2(a)=a^2+1+a\frac{\phi(a)}{\Phi(a)},
\]
the explicit marginal consequence
\begin{align}
 \|\mu_{\sigma,a}-p_d\|_1
 \le{}&|a-d|+
 \frac{2\overline\Phi(L-a)}{\Phi(a)}
 +\sqrt{\frac2\pi}u_c\sigma M_2(a)\nonumber\\
 &+2\{\overline\Phi(m_\rho/\sigma)
       +\overline\Phi((1-r)/\sigma)\}
 +\frac{2\overline\Phi(\rho/\sigma-a)}{\Phi(a)}.
 \label{eq:rh324-prefix-bound}
\end{align}
Thus $a_\sigma\to d$ implies ordinary $L^1$ convergence.  The next lemma
supplies the weighted convergence that \eqref{eq:rh324-prefix-bound} alone
does not imply.

\begin{lemma}[Uniform fourth moment of the actual prefix]
\label{lem:fourth-moment}
For every compact $\mathcal K\subset[0,\infty)$ there are
$\sigma_0,C_{\mathcal K}>0$ such that
\begin{equation}\label{eq:uniform-fourth}
 \sup_{0<\sigma<\sigma_0}\ \sup_{a\in\mathcal K}
 \int_{I_\sigma}|u|^4\,\mu_{\sigma,a}(du)
 \le C_{\mathcal K}.
\end{equation}
Consequently, whenever $a_\sigma\to d$,
\begin{equation}\label{eq:sector-moment-convergence}
 \int_{\{\pm u>0\}}u^2\,\mu_{\sigma,a_\sigma}(du)
 \longrightarrow
 \int_{\{\pm u>0\}}u^2p_d(u)\,du.
\end{equation}
\end{lemma}

\begin{proof}
Fix $0<\rho<1-b$ and put
\[
 q_0=\max\{r,1-r\},
 \qquad z_0=\Phi(1)-\tfrac12>0.
\]
The explicit density \eqref{eq:finite-seed} gives, uniformly for
$a\in\mathcal K$ and small $\sigma$,
\begin{equation}\label{eq:V8}
 \mathbb E_{\widetilde h_{\sigma,a}}V^8\le C_{\mathcal K}.
\end{equation}
Indeed, its denominator is bounded below, the first numerator term is a
Gaussian with bounded center, and the reflected term begins at
$L-\sup\mathcal K$ after reflection and is Gaussian-tail small.

For every $x\in[0,1]$ and small $\sigma$, the folded normalizer is at least
$z_0$: the corresponding standard-normal interval has length at least two
and contains zero.  On $v\le\rho/\sigma$, the source $1-\sigma v$ remains
in the negative endpoint branch and its image has modulus at least
$m_\rho$.  Integrating the two numerator lobes in \eqref{eq:first-row}
therefore gives
\begin{align}
 \mathbb E(|U|^4\mid V=v)
 &\le z_0^{-1}\bigl\{c_\sigma(v)^4+6c_\sigma(v)^2+3\nonumber\\
 &\hspace{5.7em}{}+q_0^4\sigma^{-4}
 \overline\Phi(m_\rho/\sigma)\bigr\}.
 \label{eq:conditional-fourth}
\end{align}
Here the first three terms are the fourth moment of the full main Gaussian;
the last term bounds the remote folded lobe on the physical support.
Moreover,
\begin{equation}\label{eq:c-fourth}
 |c_\sigma(v)|^4
 \le8\alpha^4v^4+8u_c^4\sigma^4v^8.
\end{equation}
Equations \eqref{eq:V8}--\eqref{eq:c-fourth} make the contribution of
$v\le\rho/\sigma$ uniformly bounded.

On the complementary entrance tail, $|U|\le q_0/\sigma$ and Markov's
inequality applied to \eqref{eq:V8} gives
\[
 \Pr\{V>\rho/\sigma\}
 \le \rho^{-8}\sigma^8\mathbb EV^8.
\]
Its contribution to $\mathbb E|U|^4$ is therefore $O(\sigma^4)$.  This
proves \eqref{eq:uniform-fourth}.

For the second assertion, truncate $u^2\1_{\{\pm u>0\}}$ at $|u|\le M$.
The error between the truncated $\mu_{\sigma,a_\sigma}$ and $p_d$ integrals
is at most $M^2\|\mu_{\sigma,a_\sigma}-p_d\|_1$.  The two omitted tails are
at most $M^{-2}$ times their fourth moments; $p_d$ has all moments by
\eqref{eq:affine-U}.  First let $\sigma\downarrow0$ using
\eqref{eq:rh324-prefix-bound}, then let $M\to\infty$.  This proves
\eqref{eq:sector-moment-convergence} separately for both signs.
\end{proof}

\section{Sharp transported second-hybrid row term}

The data type of the main result is fixed before measuring its error.  The
two retained $(u,w)$ path laws
\begin{equation}\label{eq:hybrid-laws}
 \mathsf H^{\rm phys}_{\sigma,a}(du,dw)
 =\mu_{\sigma,a}(du)L_\sigma(u,w)\,dw,
 \qquad
 \mathsf H^{\rm tan}_{\sigma,a}(du,dw)
 =\mu_{\sigma,a}(du)A_u(w)\,dw
\end{equation}
have the \emph{same actual physical first-leg prefix}.  Only the second row
changes.  This is precisely the second hybrid row term in the retained-path
Duhamel telescope of RH-325.

For the two repelling orientations define
\begin{equation}\label{eq:D-sector}
 D^{\pm}_{\sigma,a}
 =\int_{\{\pm u>0\}}
 \|L_\sigma(u,\cdot)-A_u\|_1\,\mu_{\sigma,a}(du),
 \qquad
 D_{\sigma,a}=D^+_{\sigma,a}+D^-_{\sigma,a}.
\end{equation}
Because $u$ is retained, $D^{\pm}_{\sigma,a}$ is exactly the unhalved
$L^1$ norm of the difference in \eqref{eq:hybrid-laws} restricted to the
corresponding sector.

\begin{lemma}[The physical boundary is negligible under the actual prefix]
\label{lem:averaged-boundary}
Uniformly for $a$ in a fixed compact subset of $[0,\infty)$,
\begin{equation}\label{eq:averaged-boundary}
 \int B_\sigma(u)\,\mu_{\sigma,a}(du)=o(\sigma).
\end{equation}
\end{lemma}

\begin{proof}
Choose
$0<\varepsilon<\min\{r,b-r\}$ and set
\[
 m_\varepsilon
 =\min\{f(r+\varepsilon),1-f(r-\varepsilon)\}>0.
\]
When $|u|\le\varepsilon/\sigma$, the two survival-function arguments in
\eqref{eq:boundary-identity} are exactly
$f(r+\sigma u)/\sigma$ and $(1-f(r+\sigma u))/\sigma$.  Hence
$B_\sigma(u)\le4\overline\Phi(m_\varepsilon/\sigma)$ there.  Off this local
chart, $B_\sigma\le2$, while \cref{lem:fourth-moment} gives
\[
 \mu_{\sigma,a}\{|u|>\varepsilon/\sigma\}
 \le C_{\mathcal K}\sigma^4/\varepsilon^4.
\]
Thus the integral is at most
$4\overline\Phi(m_\varepsilon/\sigma)
 +2C_{\mathcal K}\sigma^4/\varepsilon^4=o(\sigma)$.
\end{proof}

\begin{theorem}[Sharp physical repelling-return hybrid remainder]
\label{thm:transported-sharp}
Let $\sigma\downarrow0$ and let $a_\sigma\ge0$ satisfy
$a_\sigma\to d\ge0$.  Then, for each sign,
\begin{equation}\label{eq:sector-sharp}
 \boxed{
 \frac{D^{\pm}_{\sigma,a_\sigma}}{\sigma}
 \longrightarrow
 \sqrt{\frac2\pi}\,u_c
 \int_{\{\pm u>0\}}u^2p_d(u)\,du>0.}
\end{equation}
Summing the orientations gives the closed formula
\begin{equation}\label{eq:total-sharp}
 \boxed{
 \frac{D_{\sigma,a_\sigma}}\sigma
 \longrightarrow
 \sqrt{\frac2\pi}\,u_c
 \left[
 1+\alpha^2\left(
 d^2+1+d\frac{\phi(d)}{\Phi(d)}
 \right)
 \right]>0.}
\end{equation}
\end{theorem}

\begin{proof}
Let $E_\sigma(u)=\|L_\sigma(u,\cdot)-A_u\|_1$ and let
$S_\sigma(u)$ be \eqref{eq:shift-identity}.  The reverse and ordinary
triangle inequalities give
$|E_\sigma(u)-S_\sigma(u)|\le B_\sigma(u)$.  Therefore
\cref{lem:averaged-boundary} shows, in either sector,
\begin{equation}\label{eq:E-to-S}
 \left|\int_{\{\pm u>0\}}(E_\sigma-S_\sigma)
 \,d\mu_{\sigma,a_\sigma}\right|=o(\sigma).
\end{equation}

The elementary Gaussian-shift inequality
\[
 0\le4\Phi(t/2)-2\le\sqrt{\frac2\pi}\,t,
 \qquad t\ge0,
\]
implies
\begin{equation}\label{eq:curvature-domination}
 0\le\frac{S_\sigma(u)}\sigma
 \le\sqrt{\frac2\pi}u_cu^2,
 \qquad
 \frac{S_\sigma(u)}\sigma
 \longrightarrow\sqrt{\frac2\pi}u_cu^2.
\end{equation}
The convergence is uniform on each fixed bounded $u$-interval.  On such an
interval, \eqref{eq:rh324-prefix-bound} transfers the integral from
$\mu_{\sigma,a_\sigma}$ to $p_d$.  Outside it,
\eqref{eq:curvature-domination} and the uniform fourth moment make the
second-moment tails uniformly small.  Thus
\[
 \int_{\{\pm u>0\}}\frac{S_\sigma(u)}\sigma
 \,\mu_{\sigma,a_\sigma}(du)
 \longrightarrow
 \sqrt{\frac2\pi}u_c
 \int_{\{\pm u>0\}}u^2p_d(u)\,du.
\]
Together with \eqref{eq:E-to-S}, this proves \eqref{eq:sector-sharp}.
The density \eqref{eq:pd} is strictly positive on $\R$, so both sector
coefficients are strictly positive.

Finally \eqref{eq:affine-U} gives
\[
 \mathbb EU^2=1+\alpha^2\mathbb EV^2,
 \qquad
 \mathbb EV^2=d^2+1+d\frac{\phi(d)}{\Phi(d)}.
\]
Adding the two sector moments proves \eqref{eq:total-sharp}.
\end{proof}

\begin{corollary}[Scoped negative and fixed first-alias phase]
\label{cor:negative-phase}
For every $d\ge0$ and along every sequence $a_\sigma\to d$, the exact
second hybrid row error in \eqref{eq:D-sector} is neither exponentially
small nor $o(\sigma)$.
If the inherited first-alias clearance satisfies
\[
 \delta_k=C_{\rm b}\lambda^{-2k}\{1+o(1)\},
 \qquad C_{\rm b}>0,
\]
and
\[
 a_\sigma=\frac{\delta_k}{\sigma},
 \qquad
 \eta_\sigma:=k-\frac{\log(1/\sigma)}{2\log\lambda}
 \longrightarrow\eta,
\]
then
$a_\sigma\to d=C_{\rm b}\lambda^{-2\eta}$ and
\eqref{eq:sector-sharp}--\eqref{eq:total-sharp} hold with this phase-matched
$d$.
\end{corollary}

The comparison $\sigma=o(R^{-2k})$ for $R=1.4$ and
$k=\log(1/\sigma)/(2\log\lambda)+O(1)$ remains valid because
$\log R/\log\lambda=0.6496301165\ldots<1$.  This is only local scale
compatibility.  The positive coefficient is an obstruction to improving
the hybrid term below order $\sigma$, not an obstruction at the larger
first-alias target scale.

\section{Deterministic reproduction and claim boundary}

The implementation evaluates the exact normalizer, both row identities,
triangle bounds, the density $p_d$, and its two sector moments using
standard-library composite Simpson quadrature.  The displayed sector
integrals use the finite numerical window $|u|\le32$; the exact total uses
the closed formula in \eqref{eq:total-sharp}.  Selected coefficients are
\begin{center}
\begin{tabular}{c@{\quad}rrr}
\toprule
$d$ & $u<0$ coefficient & $u>0$ coefficient & exact total\\
\midrule
$0$   & $12.88903630$ & $0.08295737$ & $12.97199367$\\
$0.5$ & $18.83879635$ & $0.05712449$ & $18.89592085$\\
$1$   & $28.05382711$ & $0.03498685$ & $28.08881396$\\
\bottomrule
\end{tabular}
\end{center}
The sector sums agree with the closed total to the displayed digits.  Finite
curvature-proxy rows approach these values and exact row quadrature lies
inside \eqref{eq:triangle-sandwich}.  These rows check formulas; no exponent
is fitted and no finite table is used to prove an asymptotic statement.

The signs $U<0$ and $U>0$ label the two orientations on either side of the
repelling coordinate $r$.  They are \emph{not} the two RH-19 critical
siblings, which live in $\sqrt\sigma$-scale windows around the distinct
partition point $b$.  The folded companion lobe in
\eqref{eq:physical-second-row} is already included in the exact boundary
identity and is not assigned a parity weight.

\begin{remark}[Exact data-type firewall]\label{rem:firewall}
The equality represented by $D_{\sigma,a}^{\pm}$ belongs only to the second
hybrid Duhamel row term with the same actual first-leg prefix and retained
$u$.  Marginalizing $u$ gives only the contraction bound
\[
 \|\mathsf H^{\rm phys}_{\sigma,a,W}
   -\mathsf H^{\rm tan}_{\sigma,a,W}\|_1
 \le D_{\sigma,a}.
\]
Replacing the first physical leg as well creates a separate hybrid term;
there is no asserted equality between the fully physical two-leg law and
the fully affine two-leg law.  Nor does the theorem supply a global
uniform row bound, an all-cycle $O(k\sigma)$ estimate, phase transport
through $2k$ legs, cyclic trace control, parity/shell cancellation,
full-trace replacement, or determinant gluing.
\end{remark}

Gates A--E remain false/open.  This paper constructs no Hilbert--Polya
operator, identifies no Riemann zero, proves no von Mangoldt trace formula
or completed-zeta divisor equality, and does not imply the Riemann
Hypothesis.  The next legitimate interface, RH-333, is to transport the
common clearance phase and incoming moments through the full boundary cycle
and either prove the retained-path $O(k\sigma)$ estimate on that one clock
or isolate a genuine phase/stability obstruction.  Its first admissible
test is the possible escape gap in the raw mass-one forward all-affine
chain.  That test must be kept distinct from any cyclic bridge, Doob
transform, or branch-complete non-Gaussian reference, which is not yet
defined in the required physical data type.

\end{document}
